\section{At the Evening of Life}\label{sec:particular}

``It is appointed unto men once to die, and after this the judgment'' (Heb 9:27). The Apostle does not say how long after, and the Church does not leave the question to conjecture. She teaches that the judgment of each soul follows death at once: ``Each man receives his eternal retribution in his immortal soul at the very moment of his death, in a particular judgment that refers his life to Christ: either entrance into the blessedness of heaven---through a purification or immediately,---or immediate and everlasting damnation.''\endnote{\emph{Catechism of the Catholic Church} 1022, official English.} The Catechism sets beside that sentence a single line of St John of the Cross, which has become the briefest summary of the whole matter: ``At the evening of life, we shall be judged on our love.''\endnote{John of the Cross, \emph{Dichos} 64, as quoted at \emph{Catechism of the Catholic Church} 1022.}

The Scriptures speak of this judgment first by its effects. The beggar Lazarus dies and is ``carried by the angels into Abraham's bosom,'' while the rich man is buried and lifts up his eyes in torment, and between them ``there is fixed a great chaos'' that no one can cross (Luke 16:22--26). To the thief who asks to be remembered Christ answers: ``Amen I say to thee: This day thou shalt be with me in paradise'' (Luke 23:43). Paul, weighing life against death, is ``straitened between two: having a desire to be dissolved and to be with Christ, a thing by far the better'' (Phil 1:23); and he writes to the Corinthians that ``we are confident and have a good will to be absent rather from the body and to be present with the Lord'' (2~Cor 5:8). Aquinas rests the whole doctrine on that last verse. Against those who held that the souls of the saints would not be rewarded, nor the souls of the lost punished, until the judgment day, he answers that to be present with the Lord is no longer to walk by faith but by sight, and that this is ``to see God in His Essence, wherein consists `eternal life'\,''; and he concludes: ``Hence it is manifest that the souls separated from bodies are in eternal life.''\endnote{Thomas Aquinas, \emph{Summa theologiae} III, q.~59, a.~5, ad~1, trans.\ Fathers of the English Dominican Province (Project Gutenberg edition of the Third Part).}

\subsection{\emph{Benedictus Deus}}

What the schools commonly taught, the Church defined. In the 1330s a dispute had arisen in which the Pope himself had taken part: whether the purified souls of the just, separated from their bodies, see the divine essence with the vision the Apostle calls face to face before they take up their bodies again. John~XXII closed his own share in it with a declaration of 3~December 1334, confessing that the purified souls separated from their bodies are in heaven with Christ in the company of the angels, and see God and the divine essence face to face, clearly, so far as the condition of a separated soul allows: \latin{vident Deum de communi lege ac divinam essentiam facie ad faciem clare, in quantum status et condicio compatitur animae separatae}.\endnote{John~XXII, bull \emph{Ne super his} (3 December 1334), DS 990--991.} On 29~January 1336 his successor Benedict~XII settled the matter for the whole Church in the constitution \latin{Benedictus Deus}, and he did so in the most solemn form a pope can use: \latin{Hac in perpetuum valitura Constitutione auctoritate Apostolica diffinimus}: by a constitution to stand for ever, and by apostolic authority, he defines.\endnote{Benedict~XII, constitution \emph{Benedictus Deus} (29 January 1336), DS 1000. Latin of DS 1000--1002 from the \emph{Enchiridion symbolorum} as presented at patristica.net, \url{https://patristica.net/denzinger/enchiridion-symbolorum.html} (accessed 16 October 2026).}

The definition is worth reading phrase by phrase, because every phrase carries weight.

It speaks first of whom it defines: the souls of the saints who died before Christ's Passion; of the apostles, martyrs, confessors, virgins, and all the faithful who died after receiving his baptism; and of baptized children who die before the use of free will. It then distinguishes, among the adult faithful, two conditions. There are those \latin{in quibus nihil purgabile fuit, quando decesserunt}---in whom there was nothing to be purged when they died; and there are those in whom there was something to be purged, \latin{cum post mortem suam fuerint purgatae}, when they have been purged after their death. Purgatory enters the definition here, as a clause within the definition of the blessed. It is not set over against heaven as a third destiny; it is the way into heaven for those who need it.

Then comes the word on which the whole constitution turns: \latin{mox}. Straightway after their death, and after that purification in those who needed it, \latin{etiam ante resumptionem suorum corporum et iudicium generale}---even before the resumption of their bodies and the general judgment---these souls ``have been, are and will be in heaven, in the heavenly Kingdom and celestial paradise with Christ, joined to the company of the holy angels.'' And there they have seen and do see the divine essence \latin{visione intuitiva et etiam faciali}, ``with an intuitive vision, and even face to face, without the mediation of any creature,''\endnote{\emph{Benedictus Deus}, DS 1000; the English of the central clauses is that of \emph{Catechism of the Catholic Church} 1023, official English.} the divine essence showing itself to them \latin{immediate \ldots\ nude, clare et aperte}, immediately, nakedly, clearly and openly. Seeing, they enjoy; enjoying, they are truly blessed and have eternal life and rest. The vision once begun is never interrupted, but continues \latin{usque ad finale iudicium et ex tunc usque in sempiternum}, until the final judgment and from then on for ever.\endnote{\emph{Benedictus Deus}, DS 1001.}

The constitution then turns, in a single sentence, to the other side: \latin{animae decedentium in actuali peccato mortali mox post mortem suam ad inferna descendunt}---the souls of those who die in actual mortal sin go down to hell straightway after their death, to be punished there.\endnote{\emph{Benedictus Deus}, DS 1002.} The same \latin{mox} stands on both sides. There is no waiting-room between the deathbed and the soul's destiny, no sleep of the soul, no interval in which a life already ended is still under deliberation. The particular judgment is not a preliminary hearing. It is the sentence itself, passed and carried out at the instant of death.

\subsection{What the moment holds}

The saints who taught the faithful to prepare for death meditated on that instant with great concreteness. Alphonsus, opening his consideration on the particular judgment with Paul's ``We must all be manifested before the judgment-seat of Christ,'' asks the reader to consider ``the appearance, the accusation, the examination, and the sentence.'' Of the first he writes: ``it is the common opinion of theologians, that the particular judgment takes place at the very moment of death; and that on the very spot where the soul is separated from the body, she is judged by Jesus Christ, who shall not send, but will come himself to judge her according to her works.''\endnote{Alphonsus Liguori, \emph{Preparation for Death}, Consideration~XXIV, first point (Duffy, 1843), p.~289.} He does not spare the reader the fear of it, and he cites Augustine to say that Christ will come ``in love to the good, in terror to the wicked.'' But his first response is not to describe the terror further. It is to pray:

\begin{quote}
O my Jesus, I will always call you Jesus. Your name consoles and encourages me, because it reminds me that you are my Saviour, who have died for my salvation. \ldots\ ``Recordare, Jesu pie, quod sum causa tu\ae\ vi\ae.'' Pardon me, then, immediately, O my Jesus, before you come to judge me. I shall not then be able to ask pardon; I can now ask it from you, and I hope for it. Your wounds shall then fill me with terror, but now they give me confidence.\endnote{Alphonsus Liguori, \emph{Preparation for Death}, Consideration~XXIV, first point, Affections and Prayers (Duffy, 1843), p.~292.}
\end{quote}

The line he borrows is from the \latin{Dies irae}, which Irons renders: ``Think, good Jesu, my salvation / Caused Thy wondrous Incarnation.'' The saint understood that the particular judgment is not a stranger's verdict. The Judge who will meet the soul at death is the same Jesus who came down for it, sought it, and died for it, and who now waits for it to ask what at death it will be too late to ask. The whole spiritual meaning of the particular judgment lies in that word \emph{now}. Benedict~XVI says it plainly: ``With death, our life-choice becomes definitive---our life stands before the judge.''\endnote{\emph{Spe salvi} 45, official English.} Whatever a man has made of himself by the time he dies, that he is before Christ, and nothing more can be added to it or taken from it. ``In such state as a man departeth out of this life,'' Gregory the Great teaches, ``in the same he is presented in judgment before God.''\endnote{Gregory the Great, \emph{Dialogues} IV.39, in \emph{The Dialogues of Saint Gregory}, trans.\ P.~W. (Paris, 1608), ed.\ Edmund G. Gardner (London: Philip Lee Warner, 1911), p.~233, \url{https://www.tertullian.org/fathers/gregory_04_dialogues_book4.htm} (accessed 16 October 2026).}

From this the masters of the spiritual life drew one practical conclusion above all others: to be judged now, by one's own choice, rather than then. Alphonsus makes his own the saying he attributes to Bernard, \latin{Volo judicatus praesentari, non judicandus}---I wish to be presented already judged, not to be judged---and turns it into a prayer: ``O my Judge, I wish to be judged and punished during life, which is a time of mercy and pardon; for, after death shall be the time of justice.''\endnote{Alphonsus Liguori, \emph{Preparation for Death}, Consideration~XXIV, third point (Duffy, 1843), p.~299.} Every confession made in sorrow is such a judgment. The sinner stands before Christ in the person of the priest, accuses himself, and receives a sentence of mercy; and the soul that has learned to come to that tribunal often will not come to the other as a stranger.
